% Related Work — draft by M3 (2026-08-03). Positioning against post-generation
% uncertainty (25) and training-based mitigation (26, 27) is unchanged by the
% natural-speech results; only the Introduction's empirical claims were rewritten.

\section{Related Work}
\label{sec:related_work}

The reliability of large audio-language models (LALMs) has emerged as a critical bottleneck for their deployment in trustworthy interactive systems. Current literature is primarily divided into evaluation benchmarks for audio hallucinations and inference-time mitigation strategies.

\subsection{Benchmarks for Hallucination and Unanswerability in Audio-QA}
Existing audio-text evaluation protocols, such as MMAU \cite{sakshi2024mmau} and Dynamic-SUPERB \cite{huang2024dynamicsuperb}, implicitly operate under a closed-world assumption where every presented question is guaranteed to be grounded in the acoustic content. To challenge this assumption, several recent benchmarks have targeted hallucinations in audio large language models. Notably, Kuan and Lee introduced AQUA-Bench \cite{kuan2026aquabench}, which systematically evaluates audio question unanswerability across three distinct simulated settings: absent answer options, incompatible answer sets, and incompatible audio-question pairs. Similarly, Zhao et al. proposed HalluAudio \cite{zhao2026halluaudio}, constructing over 5,000 verified question-answering pairs to audit hallucinations across speech, music, and sound events. 

However, these benchmarks primarily adopt a Multiple-Choice Question (MCQ) format over non-verbal environmental sound events. While MCQ is highly effective for standardized evaluation, it fails to capture the nuances of open-ended, free-form text generation typical of real-world voice assistants. Furthermore, they do not address QA over complex spoken semantic content. Our work fills this gap by introducing a free-form spoken QA diagnostic set focused on speech comprehension rather than sound event classification, assessing the model's ability to express epistemic uncertainty without relying on predefined option constraints.

\subsection{Abstention and Reliability Mitigation for LALMs}
To mitigate the tendency of LALMs to produce confident but incorrect responses, researchers have explored both training-free and training-based alignment methods. Ma et al. \cite{ma2025reliable} systematically investigated multi-modal chain-of-thought (MCoT) prompting and supervised fine-tuning (SFT) on model-specific "I don't know" (IDK) datasets, introducing the Reliability Gain Index (RGI) to evaluate the trade-off between conservative over-refusal and hubris. Their findings indicate that reliability awareness acts as a "meta-ability" that can transfer across acoustic modalities (e.g., from speech to music). 

Other alignment efforts, such as the BALSa framework \cite{kuan2025balsa} and negative-sample synthesis in LISTEN \cite{kuan2025listen}, utilize contrastive alignment or synthetic negative samples to force LALMs to distinguish what they can hear from what they cannot. However, these methods either require expensive training pipelines or rely on post-generation uncertainty elicitation (such as semantic entropy over multiple sampled outputs, as explored in \cite{kuan2026uncertainty}). In contrast, our approach introduces a pre-generation probing mitigation technique. By extracting internal activation states before token generation begins, we detect unanswerability at the representation level, offering a computationally efficient alternative that avoids the generation latency of post-hoc sampling. Unlike the adapter-tuning approaches above, the LALM's own parameters remain frozen throughout: only a lightweight external probe is fitted, so the deployed model is left bit-identical to the released checkpoint.
